\section{Sua parte: abrir e definir o escopo}
\textbf{Mensagem principal:} queremos entender o que está associado à violência letal nos municípios, e os próprios dados nos obrigaram a restringir o estudo aos municípios com \nPopMin{} mil habitantes ou mais.

\subsection{O que enfatizar}
\begin{itemize}
\item A pergunta é de associação, não de causa: é um estudo observacional.
\item As quatro limitações dos dados, em linguagem simples. A última (municípios pequenos) prepara o slide do escopo.
\item No escopo, a tabela fala sozinha: nos menores, muitos zeros e oscilação grande só pelo acaso. Por isso ficamos com \nMunEsc{} municípios (\nPctPopEsc\% da população).
\item O escopo vale para \textbf{toda} a análise, exploratória e modelo.
\end{itemize}

\subsection{Perguntas prováveis}
\begin{description}[leftmargin=0pt]
\item[Por que 20 mil, e não 10 mil ou 50 mil?] É a partir desse porte que nenhum município tem taxa zero e a oscilação por acaso cai para $\pm$\nPorteOscVinte\%. Com 50 mil perderíamos boa parte do país; com 10 mil ainda há zeros.
\item[Vocês não perdem os municípios pequenos e violentos?] Sim, e isso é uma limitação declarada: as conclusões valem para municípios com \nPopMin{} mil habitantes ou mais.
\item[Por que usar homicídios registrados se há subnotificação?] É a fonte oficial e completa disponível; controlamos a qualidade do registro pela \% de mortes de intenção indeterminada.
\end{description}

\subsection{Transição}
Ao terminar o escopo: ``Com o escopo definido, o Fabio apresenta de onde vêm os dados e as variáveis.''
